\documentclass[11pt]{article}
\usepackage[a4paper,margin=1in]{geometry}
\usepackage{graphicx}
\usepackage{booktabs}
\usepackage{amsmath}
\usepackage{amssymb}
\usepackage{hyperref}
\usepackage{natbib}
\usepackage{authblk}
\usepackage{siunitx}
\usepackage{xcolor}
\bibliographystyle{unsrtnat}

\title{A Fast Approach for Structural and Evolutionary Analysis Based on Energetic Profile Protein Comparison}
\date{}

\begin{document}
\maketitle

\begin{abstract}
Fast, accurate, and interpretable protein comparison is essential to bridge the
widening gap between annotated and uncharacterized protein sequences. Structural
alignment methods such as RMSD, TM-Score, and GR-Align remain computationally
expensive, while modern deep-learning embeddings offer speed at the cost of
interpretability. We present an alignment-free descriptor that represents each
protein by a 210-dimensional vector of pairwise amino-acid interaction energies
derived from a distance-dependent knowledge-based potential. A Structural Profile
of Energy (SPE) is computed directly from atomic coordinates, and a Compositional
Profile of Energy (CPE) estimates the same descriptor from sequence composition
alone through a predictor matrix. Manhattan distance between profiles measures
dissimilarity with no alignment step. On a 251-domain CATH benchmark, CPE
reaches 100\% 1-NN accuracy in 1~second versus 67~seconds for TM-Vec and 59.2\%
to 81.5\% for classical baselines. Clustering of 143 coronavirus spike
glycoproteins yields an Adjusted Rand Index of 0.95 for CPE at tree cut 4 and 1.00
for SPE at cut 3, against 0.87 for TM-Vec, 0.49 for multiple sequence alignment,
0.73 for RMSD, and 0.56 for TM-Score. On a proteome-scale screen of 4{,}405
coronavirus models across 28 families, CPE and SPE approach 100\% accuracy and
F1. An energy-based drug separation measure correlates significantly with the
network proximity score of Cheng and colleagues on a 65-drug antihypertensive
panel. The framework unifies hierarchical classification, twilight-zone
phylogeny, proteome-scale screening, and drug-combination scoring within a
physically interpretable and scalable pairwise descriptor.
\end{abstract}

\section{Introduction}

Protein function underlies biology, medicine, and pharmacology, yet
experimental assignment of function remains slow and expensive relative to the
rate at which new proteins are sequenced. Public repositories now contain
millions of protein sequences whose functions are inferred, if at all, by
comparison to a small set of characterized proteins. Computational comparison
therefore serves as the principal lever for closing the annotation gap, and
two traditions dominate. Sequence alignment, exemplified by PSI-BLAST and its
successors~\citep{boratyn2012delta}, compares amino-acid strings against
substitution matrices and evolutionary profiles. Structural alignment,
embodied by RMSD, the TM-Score and TM-align~\citep{zhang2005tmalign}, and
graphlet-based methods such as GR-Align~\citep{maloddognin2014gralign},
compares three-dimensional coordinates and retains sensitivity at larger
evolutionary distances.

Recent advances have pushed both paradigms forward. AlphaFold2~\citep{jumper2021alphafold},
ESMFold~\citep{lin2023esmfold}, and OmegaFold~\citep{wu2022omegafold} have
eased structure prediction as a bottleneck, while learned structural embeddings
such as TM-Vec~\citep{hamamsy2024tmvec} map proteins into a vector space that
approximates TM-Score under cosine similarity. Nevertheless, alignment-based
pipelines remain expensive at proteome scale, and deep embeddings are opaque:
their dimensions do not correspond to physical quantities, and their accuracy on
rare folds depends on the distribution of training data. A representation that
was simultaneously fast, accurate on novel folds, and interpretable in terms of
protein biophysics would simplify downstream tasks from remote homology detection
to drug repositioning.

Statistical energy functions offer such a representation. Knowledge-based
potentials derived from known structures empirically capture the most probable
state of a protein and describe pairwise residue interactions within the
native fold~\citep{koppensteiner1998knowledge,sippl1990potential}. The
prevailing assumption is that a native protein occupies a minimum-energy
state: the closer a structure is to native, the lower its total energy. We
take a step beyond this and hypothesize that two proteins that are structurally
or evolutionarily similar possess analogous \emph{energy profiles}. Concretely,
each protein is summarized by a 210-dimensional vector whose entries sum
pairwise amino-acid interaction energies computed from a distance-dependent
knowledge-based potential. This Structural Profile of Energy (SPE) is computed
directly from atomic coordinates. A corresponding Compositional Profile of
Energy (CPE) is estimated from amino-acid composition alone through a
predictor matrix, eliminating the requirement for a solved structure.

The idea of using energy as a compatibility signal has deep roots. Bowie,
L\"uthy and Eisenberg~\citep{bowie1991eisenberg} pioneered environmental
energy profiles that scored the compatibility of a sequence with a structural
template. Bryngelson, Onuchic and Wolynes~\citep{bryngelson1995wolynes}
formalized the energy-landscape theory of folding via a spin-glass Hamiltonian.
Dosztanyi and colleagues~\citep{dosztanyi2008energy} estimated stabilization
energy directly from amino-acid composition. Our contribution is to collapse
such energetic information into a fixed-length pair-indexed descriptor that
supports pairwise dissimilarity via the Manhattan distance, without any
alignment step. The Manhattan distance decomposes additively over the 210
amino-acid pair types, preserving interpretability, and computes in time
linear in the descriptor length.

We validate the framework on six settings of increasing difficulty. First, we
examine whether the SPE derived from structure correlates with the CPE derived
from sequence across the ASTRAL40 and ASTRAL95 benchmark subsets of
SCOPe~\citep{chandonia2019scope}, and whether both capture the SCOP hierarchy
at class, fold, superfamily, and family levels. Second, we pit CPE and SPE
against GR-Align~\citep{maloddognin2014gralign}, RMSD, TM-Score, Yau-Hausdorff
distance, and TM-Vec~\citep{hamamsy2024tmvec} on a 251-domain CATH benchmark.
Third, we reconstruct the phylogeny of the ferritin-like
superfamily~\citep{lundin2012ferritin} beyond the twilight
zone~\citep{chothia1986twilight}, comparing against a manually curated network.
Fourth, we cluster 143 coronavirus spike glycoproteins curated from
CoV3D~\citep{gowthaman2021cov3d}. Fifth, we classify 690 bacteriocins from the
BAGEL database~\citep{alkhalili2018bagel} and four Betacoronavirus subgenera of
the Papain-like protease domain. Finally, we extend the method to drug
combinations by correlating a new energy-profile separation measure with the
network proximity score of Cheng and colleagues~\citep{zhou2020network} on a
panel of 65 antihypertensive drugs. Together with a proteome-scale screen of
4{,}405 coronavirus protein models across 28 families, these experiments
establish a Pareto frontier that is simultaneously more accurate and
substantially faster than prior alignment-based and learning-based approaches.

\section{Related Work}

\paragraph{Knowledge-based statistical potentials.}
Statistical potentials infer effective pairwise energies from the observed
frequencies of atom--atom contacts in high-resolution structures. Sippl's
seminal formulation~\citep{sippl1990potential} and later
refinements~\citep{koppensteiner1998knowledge} established distance-dependent
potentials as a workhorse for threading, fold recognition, and model quality
assessment. We adopt a distance-dependent variant of this family of potentials
as the foundation of our descriptor.

\paragraph{Energy landscape theory.}
Energy-landscape ideas, culminating in the
Bryngelson--Wolynes funnel picture~\citep{bryngelson1995wolynes}, cast the
native fold as a global energy minimum reached through minimally frustrated
pathways. Recent work has exploited residual frustration to probe functional
divergence in close homologs such as $\alpha$- and
$\beta$-globins~\citep{freiberger2023frustration}, motivating our use of
pairwise energy distributions to distinguish close subfamilies within the
hemoglobin tetramer.

\paragraph{Structural comparison metrics.}
TM-align~\citep{zhang2005tmalign} combines dynamic programming with a
length-normalized scoring function; its TM-Score has become the de-facto
standard for reporting structural similarity. GR-Align~\citep{maloddognin2014gralign}
extracts topological graphlets from residue contact graphs for flexible
alignment. RMSD after least-squares superposition is simple and widely used
but sensitive to outlier regions and flexible loops. Hausdorff-type distances
treat structures as point sets and measure maximal bidirectional separation.
These methods differ in their balance of accuracy, speed, and robustness, but
all of them share a dependence on either alignment or explicit atomic
coordinates.

\paragraph{Deep-learning embeddings.}
TM-Vec~\citep{hamamsy2024tmvec} encodes structures into a fixed-length vector
using a transformer trained on structure pairs with TM-Score supervision;
cosine similarity in the embedding approximates TM-Score. The method achieves
strong performance on remote homology detection, but inherits the classical
limitations of deep learning: reliance on the training distribution and
limited interpretability of the learned dimensions. Our energy profile
complements such embeddings by offering dimensions that correspond directly to
biophysical interaction types.

\paragraph{Structure prediction and classification.}
AlphaFold2~\citep{jumper2021alphafold}, OmegaFold~\citep{wu2022omegafold}, and
ESMFold~\citep{lin2023esmfold} have transformed the availability of protein
structure models. Hierarchical classifications such as
SCOP~\citep{murzin1995scop}, CATH~\citep{orengo1997cath}, and
SCOPe~\citep{chandonia2019scope} organize the protein universe at class, fold,
superfamily, and family levels. Phylogenetic reconstruction in the
twilight zone~\citep{chothia1986twilight}---where sequence similarity falls
below 20\%--35\%---requires structural signal; Lundin and
colleagues~\citep{lundin2012ferritin} constructed structural evolutionary
networks for the ferritin-like superfamily using methods from
phangorn~\citep{schliep2011phangorn} and SplitsTree~\citep{huson2024splitstree}.
Multiple-sequence alignments obtained via MAFFT~\citep{rozewicki2021mafft} and
non-redundant set curation via PISCES~\citep{wang2003pisces} remain standard
supporting tools. CoV3D~\citep{gowthaman2021cov3d} curates high-resolution
coronavirus structures and their antibody complexes.

\paragraph{Network pharmacology.}
Cheng and colleagues introduced a network-based separation measure $S_{AB}$
between pairs of drug targets in the human
interactome~\citep{zhou2020network}. Drug pairs with $S_{AB}<0$ share a
network neighborhood, while pairs with $S_{AB}>0$ are topologically separated.
We define an analogous energy-based separation measure $E_{AB}$ on the drug
targets' energy profiles and show that the two measures are strongly
correlated on a 65-drug antihypertensive panel, without requiring network
reconstruction.

\section{Methods}
\label{sec:methods}

\subsection{Training dataset and filtering}

The knowledge-based potential is trained on a non-redundant dataset of protein
chains curated from PISCES~\citep{wang2003pisces}. Chains satisfy pairwise
sequence identity below 50\%, crystallographic resolution higher than
$1.6~\si{\angstrom}$, R-factor below 0.25, and length between 40 and 1{,}000
residues. Chains overlapping with any downstream test set are removed. This
selection yields a diverse and high-quality training set for fitting the
potential without contaminating the evaluation benchmarks.

\subsection{Distance-dependent knowledge-based potential}

Each protein chain is represented by its heavy atoms, and contacts between
atoms are identified via Delaunay tessellation implemented through
Qhull~\citep{barber1996qhull}. Two atoms are considered to be in contact when
they are connected by an edge in the Delaunay triangulation, ensuring direct
interaction without obstruction. Interatomic distances are discretized into
$30$ shells of width $0.5~\si{\angstrom}$ starting at $0.75~\si{\angstrom}$.
Following the Sippl formalism~\citep{sippl1990potential,koppensteiner1998knowledge},
$167$ atom types are defined by treating non-hydrogen atoms as distinct
according to their amino-acid residue. The energy between atoms $i$ and $j$
at distance $d$ is
\begin{equation}
  E_{ij}(d) = RT \, \ln \left( 1 + M_{ij}\sigma \right) - RT \, \ln \left(1 + M_{ij}\sigma\, \dfrac{f_{ij}(d)}{f_{xx}(d)} \right),
  \label{eq:sippl}
\end{equation}
with $RT = 0.582~\text{kcal/mol}$, $\sigma = 0.02$, $M_{ij}$ the number of
observations for atomic pair $(i,j)$, $f_{ij}(d)$ the relative frequency of
$(i,j)$ in shell $d$, and $f_{xx}(d)$ the relative frequency for all atomic
pairs in shell $d$. The residue-residue energy $\Delta E(A,B)$ between amino
acids $A$ and $B$ is the sum of $E_{ij}(d_{ij})$ over atom pairs in contact.

\subsection{Structural Profile of Energy}

For each protein structure, a $210$-dimensional vector is built whose entries
index unordered amino-acid pairs. The entry for pair $(A,B)$ is the sum of
$\Delta E(A,B)$ over all residue pairs of types $(A,B)$ in contact. This
vector is the Structural Profile of Energy (SPE). Because the potential is
distance-dependent and residue types account for a large share of chemical
context, the SPE collapses a full residue interaction graph into a compact
descriptor that preserves pairwise interaction information.

\subsection{Compositional Profile of Energy from sequence}

Many proteins of interest lack experimentally determined three-dimensional
structure. We therefore introduce a sequence-only descriptor that estimates
the pairwise energy content from amino-acid composition. Let $n_i$ denote the
frequency of amino-acid type $i$ in a protein $S$ of total length $L$, and
$f_i = n_i / L$. We define a predictor matrix $P$ of shape $20 \times 20$
such that the estimated energy involving amino-acid type $i$ is
\begin{equation}
  \widehat{E}_i(S) = L \sum_{j=1}^{20} P_{ij} f_j,
  \label{eq:predictor}
\end{equation}
and the pairwise Compositional Profile of Energy (CPE) entry for types
$(i,j)$ is derived from $P_{ij}$ and the composition frequencies. Each row of
$P$ is fit independently by minimizing the squared error between
$\widehat{E}_i(S)$ and the structural reference $E_i(S) = \sum_{j}
\Delta E(i,j;S)$ over the PISCES training set. Setting the partial
derivatives of the objective to zero yields a linear system that we solve
with the MATLAB Symbolic Math Toolbox.

\subsection{Dissimilarity and drug-pair separation}

For a pair of proteins $(a,b)$, the dissimilarity $d(a,b)$ is the Manhattan
distance between their 210-dimensional energy profiles. For a pair of drugs
$(A,B)$ with target sets $T_A$ and $T_B$, we define an energy separation
measure
\begin{equation}
  E_{AB} = \langle d(a,b) \rangle_{a \in T_A, b \in T_B} - \tfrac{1}{2}\big(\langle d(a,a') \rangle_{a,a' \in T_A} + \langle d(b,b') \rangle_{b,b' \in T_B}\big),
  \label{eq:eab}
\end{equation}
analogous in form to the network-based separation $S_{AB}$ introduced by
Cheng and colleagues~\citep{zhou2020network}. Negative values indicate
targets that are energetically close; positive values indicate energetically
separated target modules.

\subsection{Benchmarks and evaluation}

We use several benchmarks. For hierarchical structure, ASTRAL40 and ASTRAL95
subsets of SCOPe~v2.08~\citep{chandonia2019scope} provide domains with
sequence identity below 40\% and 95\% respectively, filtered by E-value.
For head-to-head comparison, a CATH~v4.2.0 benchmark of $251$ domains drawn
from families 1.10.8.10 (Ruva Protein fold) and 3.10.28.10
(Endonuclease I-creI fold), with lengths between 44 and 854 residues
(mean 211), supports 1-NN classification. Five SCOP superfamilies
(a.4.5, b.55.1, d.17.4, d.15.1, and b.1.1) test superfamily-level accuracy.
For phylogeny, the ferritin-like superfamily a.25.1 is used with its two
curated families a.25.1.1 and a.25.1.2. For coronaviruses, we curate 143
closed-RBD spike glycoprotein chains from CoV3D~\citep{gowthaman2021cov3d}
(80 SARS-CoV-2, 31 SARS-CoV, 32 MERS-CoV), and the Freiberger dataset of
4{,}405 models spanning 28 families for large-scale screening. For
bacteriocins, we use 690 proteins from BAGEL~\citep{alkhalili2018bagel} with
length greater than $30$ amino acids. For drug combinations, we adopt the
65-drug antihypertensive panel from Cheng and colleagues. Evaluation metrics
include 1-NN accuracy, F1 score, the Adjusted Rand Index, Spearman rank
correlation, and bootstrap resampling with 100 replicates. Reported runtimes
are measured on a $2.4~\si{\GHz}$ processor with $8~\si{\GB}$ of RAM.

\begin{figure}[t]
  \centering
  \includegraphics[width=0.95\linewidth]{figures/fig_pipeline.pdf}
  \caption{Development of the knowledge-based potential and energy profiles.
    Non-redundant PISCES chains are used to fit a distance-dependent
    knowledge-based potential whose atomic contacts are identified by
    Delaunay tessellation. The Structural Profile of Energy (SPE) summarizes
    pairwise residue energies for a given structure; the Compositional
    Profile of Energy (CPE) estimates the same 210-dimensional descriptor
    from amino-acid composition using a predictor matrix $P$. Dissimilarity
    is quantified by the Manhattan distance.}
  \label{fig:pipeline}
\end{figure}

\section{Results}

\subsection{Sequence-derived energy matches structural energy}

A necessary prerequisite for a sequence-based descriptor is that it tracks
its structural counterpart. We therefore examined whether total energy
estimated from sequence via the predictor matrix agrees with total energy
computed directly from three-dimensional coordinates. On the ASTRAL40 and
ASTRAL95 datasets, we computed total energies from both structure and
sequence for every domain and assessed their concordance with two-sided
Pearson correlation.

Across both datasets, the total energies derived from sequence and structure
were strongly correlated, with a p-value below $10^{-16}$ (which sits below
the precision threshold of standard statistical computations). To probe
robustness, we examined the residuals---the difference between structural and
sequence-based total energies---as a function of protein length and found no
significant correlation, indicating that accuracy is independent of length.
We then repeated this analysis at the level of individual interaction types.
For each of the 210 amino-acid pair interactions, we computed the correlation
between the residual and protein length; 96\% of interaction types showed
$|r|<0.5$, consistent with the hypothesis that sequence-based energy
estimation is robust across interaction types.

These findings establish CPE as a reliable surrogate for SPE when the
three-dimensional structure is unavailable. Pairwise distances between
energy profiles derived from sequence and from structure were also strongly
correlated on both ASTRAL40 and ASTRAL95. Finally, the distribution of total
energy across the four SCOP classes (all-$\alpha$, all-$\beta$,
$\alpha+\beta$, $\alpha/\beta$) differed systematically and was consistent
between the sequence- and structure-based estimates, in line with the
recognized energetic differences between predominantly helical and
$\beta$-strand folds.

\subsection{Energy profiles encode the SCOP hierarchy}

A protein descriptor worth its name should expose the hierarchical structure
of the protein universe. We tested this by projecting the 210-dimensional
profiles into two dimensions with UMAP. To examine class-level structure, we
embedded all-$\alpha$ and all-$\beta$ domains from ASTRAL40 and ASTRAL95 and
found clearly distinct energy patterns. To examine finer levels, we
randomly selected two folds (a.100 and a.104) from the all-$\alpha$ class,
two superfamilies (a.29.2 and a.29.3) from fold a.29, and two families
(a.25.1.0 and a.25.1.2) from superfamily a.25.1, and projected their CPE and
SPE vectors with $n_\text{neighbors}=30$ and $\text{min\_dist}=0.1$.

The embeddings show that domains sharing the same fold, superfamily, or
family cluster together in both the CPE and SPE views. Quantitatively,
intraclass distances within the all-$\alpha$ class from ASTRAL95 were
significantly smaller than interclass distances, and the same pattern held
when pairwise distances within fold a.29 were compared to distances across
different folds within the all-$\alpha$ class, and again for pairs within
superfamily a.29.3 versus pairs within fold a.29 that belong to different
superfamilies. Hierarchical structure is therefore preserved in the energy
profile across all four SCOP levels examined.

\subsection{Head-to-head comparison on CATH and five SCOP superfamilies}

We next benchmarked the descriptor against established structural metrics on
a $251$-domain CATH subset drawn from families 1.10.8.10 and 3.10.28.10. A
1-nearest-neighbor classifier was paired with GR-Align, RMSD, TM-Score, the
Yau-Hausdorff distance, TM-Vec, SPE, and CPE as dissimilarity functions.

Both CPE and TM-Vec reach 100\% accuracy in distinguishing the two protein
families, whereas GR-Align, RMSD, TM-Score, and the Yau-Hausdorff distance
fall between 59.2\% and 81.5\%. The computational gap is pronounced. CPE
completes the full 1-NN classification in 1~second, while TM-Vec requires
67~seconds on the same $2.4~\si{\GHz}$ / $8~\si{\GB}$ RAM hardware. SPE,
which requires the coordinates but no alignment, is similarly competitive.
Figure~\ref{fig:accuracy_time} summarizes both axes. Across five SCOP
superfamilies (a.4.5, b.55.1, d.17.4, d.15.1, b.1.1) we observed the same
pattern of near-ceiling accuracy, as reported in Table~\ref{tab:tab1}.

Taken together, these results establish an accuracy--time Pareto frontier
dominated by the energy-profile descriptors. Classical alignment-based
methods lag substantially in accuracy; TM-Vec matches CPE on accuracy but
costs roughly $67\times$ more compute. The interpretability advantage of
the energy descriptor---each dimension corresponds to a specific
amino-acid-pair interaction type---is free on top of these quantitative gains.

\begin{figure}[t]
  \centering
  \includegraphics[width=0.95\linewidth]{figures/fig_accuracy_time.pdf}
  \caption{Performance and computational efficiency on the 251-domain CATH
    benchmark. (A)~1-NN classification accuracy: CPE, SPE, and TM-Vec reach
    100\% while classical alignment baselines fall in the 59.2\%--81.5\% range.
    (B)~Runtime for the two methods that attain 100\% accuracy: CPE completes in
    1~second on a $2.4~\si{\GHz}$ / $8~\si{\GB}$ RAM system, whereas TM-Vec
    requires 67~seconds.}
  \label{fig:accuracy_time}
\end{figure}

\begin{table}[t]
\centering
\caption{1-NN classification summary on five SCOP superfamilies
(a.4.5 winged helix; b.55.1 PH domain-like; d.17.4 NTF-like; d.15.1
Ubiquitin-like; b.1.1 Immunoglobulins). All three methods reach
performance levels approaching $100\%$ on both accuracy and F1; CPE
matches TM-Vec on accuracy with substantially faster computation.}
\label{tab:tab1}
\begin{tabular}{lcc}
\toprule
Method & Accuracy & F1 score \\
\midrule
CPE    & approaching 100\% & approaching 100\% \\
SPE    & approaching 100\% & approaching 100\% \\
TM-Vec & approaching 100\% & approaching 100\% \\
\bottomrule
\end{tabular}
\end{table}

\subsection{Fine-grained discrimination of $\alpha$- and $\beta$-globins}

$\alpha$- and $\beta$-globin subunits share a common evolutionary origin and
highly similar folds, but perform distinct roles within the hemoglobin
$\alpha_2 \beta_2$ tetramer. They therefore provide a stringent test for
subfamily-level resolution. Using a non-redundant set of $21$ mammalian
hemoglobin structures in which each subunit exhibits its own pattern of
conserved energetic frustration~\citep{freiberger2023frustration}, we
computed CPE, SPE, and TM-Vec embeddings and visualized them with UMAP
using $n_\text{neighbors}=13$ and $\text{min\_dist}=0.1$.

All three methods cleanly separate $\alpha$- from $\beta$-globins in the
embedding. Because the two subfamilies carry comparable global structural
scores under conventional structural metrics, this test demonstrates that
pair-indexed energetic signatures capture the subtle energetic
reorganization that accompanies functional specialization within a
homologous family. This result is consistent with the interpretation that
residual frustration patterns encode functional divergence, and suggests
that energy profiles can serve as a compact proxy for such patterns.

\subsection{Phylogeny of the ferritin-like superfamily beyond the twilight zone}

Sequence similarity in the ferritin-like superfamily often falls below the
20\%--35\% twilight zone where meaningful phylogenetic inference from
sequence alone breaks down~\citep{chothia1986twilight,lundin2012ferritin}.
Using the same structures as the manually curated reference networks, we
reconstructed phylogenetic networks from CPE, SPE, and TM-Vec pairwise
distances with the NeighborNet method, visualized them with
SplitsTree~\citep{huson2024splitstree}, and constructed neighbor-joining
trees from average inter-subfamily distances with
phangorn~\citep{schliep2011phangorn}.

All three methods correctly split the superfamily a.25.1 into its two
families a.25.1.1 and a.25.1.2. Within each family, however, the ordering
of subgroups diverges. To quantify the deviation from the manual network,
we computed Spearman rank correlations between the predicted branching
orders and the reference. The energy descriptors outperform TM-Vec on this
task (Table~\ref{tab:tab2}); SPE attains a perfect correlation on
family a.25.1.1, while TM-Vec records a negative correlation on family
a.25.1.2. The interpretation is that while TM-Vec is effective at clustering
similar proteins, it is less faithful to the finer subfamily ordering that
follows from structural evolution; in contrast, the SPE and CPE orderings
align with the manually curated reference produced from structure-based
NeighborNet analyses.

\begin{table}[t]
\centering
\caption{Spearman rank correlation between predicted and manually curated
branching orders for each subfamily of the ferritin-like superfamily.
SPE attains perfect agreement on family a.25.1.1; TM-Vec is the only method
with a negative correlation on family a.25.1.2.}
\label{tab:tab2}
\begin{tabular}{lccc}
\toprule
Family & SPE & TM-Vec & CPE \\
\midrule
a.25.1.1 (ferritin) & perfect ($+1$) & partial & strong positive \\
a.25.1.2 (RNR-like) & strong positive & negative & positive \\
\bottomrule
\end{tabular}
\end{table}

\subsection{Clustering of coronavirus spike glycoproteins}

Coronavirus spike glycoproteins have been central to structural and vaccine
research across SARS-CoV, SARS-CoV-2, and MERS-CoV outbreaks. We curated
$143$ spike chains with a closed receptor-binding domain from
CoV3D~\citep{gowthaman2021cov3d}, comprising 80 SARS-CoV-2, 31 SARS-CoV, and
32 MERS-CoV structures. We built pairwise distance matrices with each of
MSA (ClustalW via the msa package), RMSD, TM-Score, TM-Vec, CPE, and SPE,
and constructed dendrograms using the UPGMA method.

All methods separated the SARS-CoV/SARS-CoV-2 lineage from the
MERS-CoV subgenus, but differed substantially in how cleanly the three
species partitioned. Using the Adjusted Rand Index, CPE achieved the
highest sequence-level clustering at $0.95$ with a tree cut of $4$, against
$0.87$ for TM-Vec at cut $5$ and $0.49$ for MSA at cut $3$. Among
structure-based methods, SPE achieved a perfect ARI of $1.0$ at cut $3$,
versus $0.73$ for RMSD at cut $6$ and $0.56$ for TM-Score at cut $4$. A
bootstrap analysis with $100$ replicates produced confidence intervals
supporting the main branches distinguishing the three species. The runtime
gap is dramatic: CPE completes the analysis in $0.9$~seconds and SPE in
3~minutes, while TM-Vec requires $89$~seconds, MSA $72$~seconds, RMSD
70~minutes, and TM-Score $9.7$~hours.
Figure~\ref{fig:spike} displays both the ARI and runtime comparisons, and
the average distances between the three viral groups consistently indicate
that SARS-CoV-2 is closer to SARS-CoV than to MERS-CoV.

These findings make a concrete case that energy profiles recover
species-level taxonomy within a single protein family with higher fidelity
and orders-of-magnitude lower cost than either alignment-based or
deep-learning alternatives. The combination of ARI = $1.0$ for SPE and
the $0.9$-second runtime for CPE defines a new operating point on the
accuracy--speed frontier for this class of viral clustering tasks.

\begin{figure}[t]
  \centering
  \includegraphics[width=0.95\linewidth]{figures/fig_spike_ari.pdf}
  \caption{Clustering of 143 coronavirus spike glycoproteins (80 SARS-CoV-2,
    31 SARS-CoV, 32 MERS-CoV) from CoV3D. (A)~Adjusted Rand Index at the
    best cut depth per method. (B)~Runtime on the same task (log scale).}
  \label{fig:spike}
\end{figure}

\subsection{Bacteriocin classification}

Bacteriocins are ribosomally synthesized peptides produced by bacteria with
strong antibacterial activity against closely related microbial species.
The BAGEL database organizes them into three classes according to size and
stability. We analyzed $690$ bacteriocins with length greater than $30$
amino acids~\citep{alkhalili2018bagel}, computed CPE pairwise distances,
and compared them with TM-Scores derived from structures predicted by
AlphaFold2~\citep{jumper2021alphafold}, OmegaFold~\citep{wu2022omegafold},
and ESMFold~\citep{lin2023esmfold}, as well as TM-Vec.

The UMAP projection of CPE clearly partitions bacteriocins by BAGEL
annotation. In contrast, TM-Scores computed from predicted structures
failed to separate same-class from different-class pairs cleanly, whereas
TM-Vec showed some separation but with overlap. Under CPE, comparisons of
pairwise distances across groups (different class, $n=125{,}869$; same
class, $n=111{,}349$; and subclass~1 within the same class, $n=13{,}431$)
revealed statistically significant differences in all three two-tailed
Student's $t$-tests, with p-values below $10^{-16}$---the precision
threshold of standard statistical computations.

This analysis establishes that energy profiles discriminate bacteriocins at
both class and subclass level where predicted-structure TM-Scores blur
class boundaries. The result is practically important because bacteriocin
identification is traditionally hindered by sequence and structural
diversity, and because the underlying descriptor only requires sequence.

\subsection{Papain-like protease across Betacoronavirus subgenera}

The Papain-like protease (PLPro) domain catalyzes proteolysis of viral
polyproteins and interacts with the host ubiquitin and ISG15 systems. In
the Freiberger dataset, PLPro domains are divided into four subfamilies
aligned with the Betacoronavirus subgenera: Sarbecovirus ($n=31$),
Nobecovirus ($n=11$), Merbecovirus ($n=35$), and Embecovirus ($n=45$).
We embedded their CPE, SPE, and TM-Vec representations via UMAP with
$n_\text{neighbors}=13$ and $\text{min\_dist}=0.5$.

Each subgenus forms a distinct cluster under all three representations.
The consistency of CPE, SPE, and TM-Vec at this level demonstrates that the
210-dimensional signature is sufficient to resolve genus-level structure
within a single protein family, even for small subgroups such as
Nobecovirus with $n=11$. This robustness matters because PLPro is a priority
antiviral drug target and subgenus-specific motifs have direct implications
for inhibitor selectivity.

\subsection{Energy-based drug-combination scoring}

Cheng and colleagues~\citep{zhou2020network} introduced a
network-based separation measure $S_{AB}$ between pairs of drug targets in
the human protein-protein interactome; pairs with overlapping
``footprints'' exhibit stronger pharmacological relationships. We
constructed an analogous energy separation measure $E_{AB}$ by substituting
Manhattan distances between the energy profiles of target proteins into
the same functional form (Equation~\ref{eq:eab}), and evaluated the
correlation between $E_{AB}$ and $S_{AB}$ on the 65-drug antihypertensive
panel.

The two-sided Pearson correlation between $E_{AB}$ and $S_{AB}$ is strong
across the panel. This correlation indicates that energy-profile
distances capture a substantial share of the interactome-based signal
while requiring only sequence-level input rather than a curated
protein-protein interaction network. Because the computation of $E_{AB}$
avoids shortest-path searches on a million-edge interactome, it offers a
concrete computational alternative for drug-combination scoring at scale.

\subsection{Proteome-scale screening across coronavirus families}

As a final test, we applied CPE, SPE, and TM-Vec to a dataset of $4{,}405$
coronavirus protein models drawn from 28 high-quality families assembled
by Freiberger and colleagues. Candidate proteins were originally obtained
by aligning homologous sequences to the SARS-CoV-2 reference genome
MN985325 with MAFFT~\citep{rozewicki2021mafft}, clustering with CD-HIT,
generating AlphaFold2 models~\citep{jumper2021alphafold} that met stringent
quality thresholds, and grouping them into subfamilies using S3Det.

Both CPE and SPE approach $100\%$ accuracy and F1 under 1-NN
classification, as summarized in Table~\ref{tab:tab3}. Scalability is
equally decisive. On ASTRAL95 subsets ranging from 1{,}000 to 30{,}000
proteins at intervals of 5{,}000, the processing time per amino acid grows
linearly for both CPE and TM-Vec, but the CPE slope is gentler, extending
the advantage observed on smaller benchmarks to the proteome scale.
Collectively, these results establish the energy-profile descriptor as
deployable at genome and proteome scale without sacrificing classification
accuracy.

\begin{table}[t]
\centering
\caption{Large-scale 1-NN classification on the Freiberger coronavirus
dataset ($4{,}405$ protein models spanning 28 families). Both CPE and
SPE reach performance levels near 100\%, with CPE offering the faster
runtime.}
\label{tab:tab3}
\begin{tabular}{lcc}
\toprule
Method & Accuracy & F1 score \\
\midrule
CPE    & near 100\% & near 100\% \\
SPE    & near 100\% & near 100\% \\
TM-Vec & near 100\% & near 100\% \\
\bottomrule
\end{tabular}
\end{table}

\begin{figure}[t]
  \centering
  \includegraphics[width=0.8\linewidth]{figures/fig_scalability.pdf}
  \caption{Scalability on ASTRAL95 subsets of $1{,}000$ to $30{,}000$
    proteins at intervals of $5{,}000$. CPE exhibits a gentler slope in
    processing time per amino acid than TM-Vec, as reported in the source
    experimental log.}
  \label{fig:scalability}
\end{figure}

\section{Discussion}

Our central finding is that a single 210-dimensional pair-indexed vector of
pairwise amino-acid interaction energies suffices for accurate, fast, and
interpretable protein comparison across structural, functional, and
evolutionary tasks. The SPE derived from atomic coordinates and the CPE
estimated from sequence agree at p-values below $10^{-16}$ on both the
ASTRAL40 and ASTRAL95 benchmarks, and 96\% of the 210 interaction types
yield $|r|<0.5$ between the residual and protein length. This agreement
licenses the use of CPE as a sequence-only proxy for SPE whenever an
experimental or predicted structure is unavailable. The Manhattan distance
between profiles provides a pairwise dissimilarity measure that avoids
alignment altogether, and its entries decompose additively into contributions
from individual amino-acid pair types, making interpretation natural.

The accuracy--runtime gains are substantial. On the 251-domain CATH
benchmark, CPE matches TM-Vec~\citep{hamamsy2024tmvec} at 100\% accuracy
while running in 1~second versus 67~seconds, with classical alignment
methods such as GR-Align~\citep{maloddognin2014gralign},
RMSD, TM-Score~\citep{zhang2005tmalign}, and the Yau-Hausdorff distance
lagging in the 59.2\%--81.5\% range. On the 143-structure coronavirus
spike task~\citep{gowthaman2021cov3d}, CPE achieves an Adjusted Rand Index
of $0.95$ in $0.9$~seconds while SPE attains perfect ARI $=1.0$; the
same task takes $89$~seconds with TM-Vec and $9.7$~hours with TM-Score.
On the $4{,}405$-protein Freiberger dataset spanning $28$ families, CPE and
SPE approach $100\%$ accuracy and F1 under 1-NN. These numbers reposition
the accuracy--speed Pareto frontier for protein comparison.

Beyond accuracy and speed, the energy-profile framework provides
interpretability and cross-task generality. Because each dimension of the
descriptor corresponds to an amino-acid pair type, changes in the vector
can be linked to specific interaction categories (polar--polar versus
hydrophobic--hydrophobic, for example). This property contrasts with
learned embeddings such as TM-Vec, whose dimensions lack direct biophysical
meaning, and facilitates downstream analyses ranging from mutation-impact
screening to inhibitor design. The framework's generality is illustrated
by its extension to drug-combination scoring: replacing network shortest
paths with Manhattan distances between energy profiles of protein targets
yields a strong correlation with the network-based score of Cheng and
colleagues~\citep{zhou2020network} on a 65-drug antihypertensive panel
while requiring only sequence input. The same descriptor resolves the two
families of the ferritin-like superfamily beyond the twilight
zone~\citep{chothia1986twilight,lundin2012ferritin}, separates
$\alpha$- and $\beta$-globin subfamilies that share close structure but
differ in function~\citep{freiberger2023frustration}, and partitions
bacteriocins by BAGEL class where TM-Scores on predicted structures from
AlphaFold2~\citep{jumper2021alphafold}, OmegaFold~\citep{wu2022omegafold},
and ESMFold~\citep{lin2023esmfold} blur class boundaries.

Several limitations warrant caution. First, the knowledge-based potential
is trained on observed protein structures and therefore inherits any biases
present in the PISCES~\citep{wang2003pisces} non-redundant set, including
underrepresentation of rare folds and membrane proteins. Second, although
the sequence-derived CPE is length-robust on average, protein complexity
arising from folding heterogeneity or conformational dynamics may still
affect accuracy for specific cases, and disentangling these contributions
will require additional experiments. Third, the drug-combination
application correlates with a network-based reference but has not been
validated against experimental efficacy data; practical deployment will
require prospective evaluation in disease-specific settings. Fourth, the
table-level results reported here summarize ranges or aggregate values
from the source; obtaining per-subfamily, per-method cell values requires
recomputation from primary data files beyond the extraction used here.

Several avenues stand out for future work. One direction is to
\emph{reweight} the 210-dimensional profile for specific applications, for
example amplifying hydrophobic-core interactions for remote homology
detection, or surface-exposed interactions for drug-target affinity
prediction. A second direction is to combine CPE with structure-prediction
ensembles at scale: AlphaFold2, OmegaFold, and ESMFold enable SPE
computation from predicted structures for otherwise uncharacterized
proteins, while CPE provides a structure-free cross-check at a fraction of
the cost. A third direction is to extend the descriptor beyond pairwise
interactions to higher-order motifs, potentially using tensor generalizations
of the predictor matrix $P$. Finally, the energy-profile view naturally
connects to generative modeling: conditioning a sequence generator on a
target profile could guide the design of proteins with desired fold
and interaction characteristics, bridging analysis and design within a
single descriptor. In all of these directions the core feature---a
physically grounded, fast, interpretable pairwise descriptor---will remain
the common denominator.

\section{Conclusion}

We have introduced a 210-dimensional pair-indexed energy descriptor for
proteins, computable either from three-dimensional coordinates as a
Structural Profile of Energy or from amino-acid composition alone as a
Compositional Profile of Energy, together with an alignment-free Manhattan
distance-based similarity measure. Across tasks that include SCOP
hierarchy encoding, twilight-zone phylogeny, coronavirus spike clustering,
bacteriocin and PLPro classification, proteome-scale screening of
$4{,}405$ coronavirus models, and drug-combination scoring, the descriptor
is simultaneously more accurate and substantially faster than prior
alignment-based and learning-based approaches, while retaining biophysical
interpretability. The path from here runs through application-specific
reweighting, integration with large-scale structure prediction, and
extension to higher-order interaction motifs---all within the same compact,
physically grounded framework.

\bibliography{references}

\end{document}
